\chapter{连续阅读排版样章}
\label{ch:reading}
本章用于检验长篇阅读的视觉密度. 数学示例与占位段落交替出现;
明确标为占位的段落仅用于检查字距、行距和分页, 不属于译文.

\newcommand{\SampleParagraph}{%
    这是一段连续阅读用的占位文字. 阅读者沿着正文进入下一项陈述时,
    应当能够辨认条件与结论的范围, 同时保持原有的阅读节奏.
    这里重复若干长度相近的句子, 以模拟译书中对记号的解释与过渡.
    页面上的留白需要服务于段落之间的关系, 不应使每个短结论都成为
    独立的视觉中心. 本段不引入新的数学事实, 只提供稳定的文字密度.
    \par}
\section{从定义进入论证}
\SampleParagraph
\begin{definition}[实数列的极限]
    设 $(a_n)_{n \geq 1}$ 是实数列, $a$ 是实数.
    若对每个 $\varepsilon > 0$, 都存在正整数 $N$, 使得当 $n \geq N$ 时,
    有 $|a_n - a| < \varepsilon$, 则称数列 $(a_n)$
    「\termcn{收敛}(convergent)」于 $a$.
\end{definition}
\SampleParagraph
\begin{lemma}[绝对值的基本估计]
    对任意实数 $x$、$y$, 有 $|x + y| \leq |x| + |y|$.
\end{lemma}
\begin{proposition}[常数列]
    若对每个正整数 $n$ 都有 $a_n = a$, 则 $(a_n)$ 收敛于 $a$.
\end{proposition}
\begin{proof}
    给定 $\varepsilon > 0$, 取 $N = 1$. 当 $n \geq N$ 时,
    由常数列的定义, 得 $|a_n - a| = 0 < \varepsilon$.
    因此满足收敛的定义.
\end{proof}
\SampleParagraph
\begin{theorem}[两个收敛实数列的和仍然收敛, 且其极限等于两个极限的和]
    \label{thm:sum-limit}
    设 $(a_n)$、$(b_n)$ 是实数列, 分别收敛于 $a$、$b$.
    则下列陈述成立:
    \begin{enumerate}
        \item
            数列 $(a_n + b_n)$ 收敛.
        \item
            对任意 $\varepsilon > 0$, 存在正整数 $N$, 当 $n \geq N$ 时,
            有 $|(a_n + b_n) - (a + b)| < \varepsilon$.
    \end{enumerate}
    换言之, 两个数列的极限可与有限加法交换:
    \[
        \lim_{n \to \infty} (a_n + b_n) = a + b.
    \]
\end{theorem}
\begin{proof}
    给定 $\varepsilon > 0$. 由 $(a_n)$ 收敛于 $a$, 存在正整数 $N_1$,
    当 $n \geq N_1$ 时, 有 $|a_n - a| < \varepsilon / 2$.
    同理由 $(b_n)$ 的收敛性, 存在正整数 $N_2$, 当 $n \geq N_2$ 时,
    有 $|b_n - b| < \varepsilon / 2$.

    取 $N = \max\{N_1, N_2\}$. 当 $n \geq N$ 时, 两个估计同时成立.
    运用绝对值的三角不等式, 得
    \begin{align*}
        |(a_n + b_n) - (a + b)|
        & = |(a_n - a) + (b_n - b)| \\
        & \leq |a_n - a| + |b_n - b| \\
        & < \frac{\varepsilon}{2} + \frac{\varepsilon}{2}
          = \varepsilon.
    \end{align*}
    因此 $(a_n + b_n)$ 满足以 $a + b$ 为极限的收敛条件.
\end{proof}
\SampleParagraph
\begin{corollary}[有限和]
    若 $m$ 是正整数, 且对 $1 \leq j \leq m$, 数列 $(a_{j,n})_n$
    收敛于 $a_j$, 则
    \[
        \lim_{n \to \infty} \sum_{j = 1}^{m} a_{j,n}
        = \sum_{j = 1}^{m} a_j.
    \]
    这里 $m$ 固定, 因而可反复应用定理~\ref{thm:sum-limit}.
\end{corollary}
\SampleParagraph
\begin{example}[两个简单数列]
    令 $a_n = 1 / n$、$b_n = 2 / n$. 两个数列均收敛于 $0$.
    由和的极限定理, 数列 $(3 / n)$ 也收敛于 $0$.
\end{example}
\begin{exercise}[按定义验证]
    给定 $\varepsilon > 0$, 写出使得 $n \geq N$ 时
    $3 / n < \varepsilon$ 的一个正整数 $N$.
\end{exercise}
\begin{quotationbox}
    此处为引用排版占位文字. 引用保持较小字号与楷体,
    末尾单独放置署名, 不使用证明结束符号.
    \quoteauthor{占位作者}{占位著作}
\end{quotationbox}
\section{连续短陈述与多段文字}
\SampleParagraph
\begin{lemma}[非负性]
    对任意实数 $x$, 有 $|x| \geq 0$.
\end{lemma}
\begin{lemma}[对称性]
    对任意实数 $x$, 有 $|-x| = |x|$.
\end{lemma}
\begin{proposition}[零点]
    对任意实数 $x$, $|x| = 0$ 当且仅当 $x = 0$.
\end{proposition}
\SampleParagraph
\SampleParagraph
\SampleParagraph
\ifdefined\PaginationStress
    \section{专用分页压力测试}
    以下内容故意在页底开始环境, 并使用重复占位文字.
    \newpage
    \vspace*{0.83\textheight}
    \begin{theorem}[页底题头及跨页陈述的排版占位测试]
        \label{thm:stress}
        \SampleParagraph
        \SampleParagraph
        \SampleParagraph
        \SampleParagraph
        \SampleParagraph
        \SampleParagraph
        \SampleParagraph
        \SampleParagraph
        \SampleParagraph
        \SampleParagraph
        \SampleParagraph
        \SampleParagraph
    \end{theorem}
    \begin{theorem}[以列表开始的陈述]
        \begin{enumerate}
            \item
                第一项占位文字.
            \item
                第二项占位文字.
        \end{enumerate}
    \end{theorem}
    \begin{theorem}[以公式开始的陈述]
        \[
            1 + 1 = 2.
        \]
        这是一项公式起首的排版测试.
    \end{theorem}
    对跨页陈述的引用为定理~\ref{thm:stress}.
\fi
